\section{Aula 14}
Essa aula não será possível de repassar para $\LaTeX$
sem um sacrifício considerável - talvez humano. Usamos
o teorema da aula passada [\ref{trm:lss}]
para provar resultados em $\Z^2$ orientado.

Definimos um modelo $\vec{\Z}^2 = (\Z^2, \cE)$ onde
$\cE = \{(x, x + e_i) : x \in \Z^2, 1 \leq i \leq 2\}$. Note
que as arestas desse grafo são orientadas, só podemos andar
para cima e para direita. Uma sugestão do Augusto é olhar para
esse modelo rotacionado em $45^\circ$, é um pouco mais natural.
Estamos interessados no evento $\{0 \to \infty\}$ em que existe um
caminho infinito direcionado saindo do $0$.  Definimos
como sempre
\[
	\vec{p}_c = \inf \{p \in [0,1] : \Pr_p(0 \to \infty) > 0\}.
\]

Relatamos algumas coisinhas simples, se $p$ for pequeno o suficente então
por contagem de caminhos $\Pr_p(0 \to \partial B_n) \leq \exp(-c(p)n)$.
Da mesma forma, se $p$ for grande o suficiente, podemos olhar para o dual,
fazer uma continha um pouco mais envolvida e mostrar que
$\Pr_p(0 \not \to \infty) \leq \exp(-c(p)n)$. Essas foram só bounds
pertubativos do modelo. Depois usamos LSS duas vezes para conseguir
condições finitas de percolação.

Olhando para o modelo rotacionado agora. Dado um retângulo em pé $R$ de
largura $w$ e altura $h$, definimos os eventos $V(R)$ em
que cruzamos verticalmente (de baixo para cima) o rtângulo $R$
e o evento $H(R)$ em que cruzamos horizontalmente. Aqui cruzar requer
a hipótese extra que o caminho aberto seja direcionado. Então, por exemplo,
um cruzamento horizontal sempre termina acima de onde começou.

\begin{lemma}
	Existe $\eps > 0$ tal que para todo $p$, se exitem $w,h$ tal que
	\[
		\Pr_p(V(R_{w,3h})) > 1 - \eps \quad \& \quad \Pr_p(H(R_{3w,h})) > 1 - \eps
	\]
	então $p > \vec{p}_c$.
\end{lemma}

\begin{remark}
	Se cruzar na direção difícil acontece com probabilidade alta, então percolamos.
\end{remark}

\begin{proof}
	A ideia é considerar um modelo normalizado fazendo $L$'s
	de cabeça para baixo com esses retângulos,
	é fácil ver que o modelo é algo como $3$-independente. Usando LSS
	[\ref{trm:lss}] e pedindo para
	que $\phi(1-\eps)$ seja grande o suficiente, percolamos pelo teorema
	pertubativo
	e portanto percolamos no modelo original. Além do mais, conseguimos
	de graça decaimento exponencial.
\end{proof}

Da mesma forma, olhando para o dual, conseguimos um outro lema de
condição finita.
\begin{lemma}
	Existe $\eps > 0$ tal que para todo $p$, se existem $w,h$ tal que
	\[
		\Pr_p(V(R_{3w,h})) < \eps \quad \& \quad \Pr_p(H(R_{w,3h})) < \eps
	\]
	então $p < \vec{p}_c$.
\end{lemma}
\begin{remark}
	Se cruzar na direção fácil é suficientemente difícil, então não percolamos.
\end{remark}
\begin{proof}
	A ideia aqui é fazer quase que a mesma coisa só que um grafo dual esquisito.
	Construímos uma renormalização onde o sítios são cobertos por $\Pi$
	(montado por retângulos) no dual.
	Temos $k$-independência também e usando LSS com o resultado
	pertubativo, provamos
	esse lemma.
\end{proof}

\begin{obs}
	As provas foram bem geométricas, precisaria de muito Tikz, paciência
	e parcimonia. Mas deu
	para sentir um gostinho do poder do LSS.
\end{obs}
